% !TeX root = ../main.tex

\sysusetup{
  keywords  = {推荐系统, 大语言模型, 知识图谱, 用户兴趣建模, 对比学习},
  keywords* = {recommender system, large language model, knowledge graph, user interest modeling, contrastive learning},
}

\begin{abstract}
  推荐系统是缓解信息过载的关键技术，但传统方法受限于封闭的训练数据，难以利用交互之外的世界知识。大语言模型与知识图谱分别为推荐提供了语义知识与结构化知识，然而二者的表示形式存在根本差异，简单的特征拼接难以发挥其互补优势，且大语言模型的推理开销与在线推荐的低时延要求之间存在矛盾。

  针对上述问题，本文分别从兴趣建模与知识图谱两个落点研究大语言模型在推荐中的不同角色：前者将大语言模型作为知识编码器，后者作为图构建器。其一，提出大语言模型增强的用户兴趣聚合表征推荐方法：利用大语言模型生成物品静态语义画像与用户、物品动态画像，对物品画像向量聚类得到兴趣原型中心，以目标物品的原型为查询，通过协同过滤特定注意力、物品动态注意力与用户动态注意力聚合用户兴趣表征，并引入语义难负例对比学习对齐协同嵌入与语义表示；通过离线预计算与在线内积，在保持精度的同时显著降低推理开销。其二，提出以大语言模型为图构建器的多智能体图增强推荐方法：由知识融合智能体与语义连接智能体协同生成知识连接与语义连接，与交互连接共同构成多视角图；由长尾增强智能体对稀疏用户补全伪链路、由头部去噪智能体筛除噪声交互，并在三个公开数据集上验证了其有效性。

  在多个公开数据集上的实验表明，所提方法在推荐精度与推理效率上均优于现有基线。

% TODO: 摘要定稿时压缩到 300–500 字，并把具体数值（AUC 提升幅度、加速倍数）补进去。
\end{abstract}

\begin{abstract*}
  Recommender systems are a key technology for alleviating information overload, yet conventional methods are confined to closed training data and cannot exploit world knowledge beyond user--item interactions. Large language models (LLMs) and knowledge graphs (KGs) provide semantic and structural knowledge, respectively. However, the two knowledge sources differ fundamentally in representation, and simple feature concatenation fails to exploit their complementarity; moreover, the inference cost of LLMs conflicts with the low-latency requirement of online recommendation.

  To address these issues, this thesis studies knowledge-enhanced recommendation from two complementary perspectives: vector-level and graph-structure-level knowledge injection. First, an LLM-enhanced recommendation method with aggregated user interest representation is proposed. LLMs are used to generate static semantic profiles of items and dynamic profiles of users and items; item profile vectors are clustered into interest prototypes; with the prototype of the target item as the query, user interest representations are aggregated through collaborative-filtering-specific attention, item dynamic attention, and user dynamic attention; semantic hard-negative contrastive learning aligns collaborative embeddings with semantic representations. Offline precomputation and online inner product keep inference cost low while preserving accuracy. Second, an LLM-driven multi-agent method for graph construction and graph sampling is proposed. Knowledge-fusion agents and semantic-connection agents jointly generate knowledge links and semantic links, which together with interaction links form a multi-view graph. A sampling-optimization agent completes pseudo-links for long-tail nodes and filters strong links for head nodes. Multi-view graph convolution fuses the representations of the three views to produce the final recommendation. Experiments on three public datasets verify the effectiveness of the proposed method.

  Experiments on multiple public datasets show that the proposed methods outperform state-of-the-art baselines in both accuracy and efficiency.

% TODO: 与中文摘要保持一致，定稿时校对术语。
\end{abstract*}
